% Exercise bank — Chapter 2
% The chapter is determined by this file's name.

\begin{exo}[From cold ice to boiling water: orders of magnitude]{chauffage-glace-eau-ebullition}
\seoready{true}
\keywords{calorimetry, phase change, melting, vaporization, latent heat, order of magnitude}

\textbf{Numerical data (at atmospheric pressure):}
\[
\begin{aligned}
c_{\text{ice}} &= 2{,}1\times10^3\,\text{J}\!\cdot\!\text{kg}^{-1}\!\cdot\!\text{K}^{-1},
& L_f &= 3{,}34\times10^5\,\text{J}\!\cdot\!\text{kg}^{-1},\\
c_{\text{water}} &= 4{,}18\times10^3\,\text{J}\!\cdot\!\text{kg}^{-1}\!\cdot\!\text{K}^{-1},
& L_v &= 2{,}26\times10^6\,\text{J}\!\cdot\!\text{kg}^{-1}.
\end{aligned}
\]
At atmospheric pressure, one kilogram of ice, initially at $-100\,{}^\circ\text{C}$, is heated progressively.

\begin{enumerate}
  \item Calculate the energy required to heat the ice from $-100\,{}^\circ\text{C}$ to $0\,{}^\circ\text{C}$.
  \item Calculate the energy required to melt all the ice at $0\,{}^\circ\text{C}$.
  \item Calculate the energy required to heat the liquid water from $0\,{}^\circ\text{C}$ to $100\,{}^\circ\text{C}$.
  \item Calculate the additional energy required to vaporize all the water at $100\,{}^\circ\text{C}$.
  \item Which stage requires the most energy? Compare the energy needed to heat one kilogram of water from $0\,{}^\circ\text{C}$ to $100\,{}^\circ\text{C}$ with the gravitational potential energy of a mass $m$ falling through ten metres. What mass must fall to release the same energy?
  \item Now consider the reverse process. A mass $m_v$ of water vapour condenses on a windscreen, without the resulting liquid water subsequently cooling. Express the heat $Q_{\text{rel}}$ released during condensation, then calculate it for $m_v=1{,}0\times10^{-2}\,\text{kg}$. At the windscreen temperature, take $L_v=2{,}45\times10^6\,\text{J}\!\cdot\!\text{kg}^{-1}$.
\end{enumerate}

\begin{indication}
To heat a phase without a phase change, use $Q=mc\Delta T$. For a phase change at constant temperature, use $Q=mL$. The gravitational potential energy of a mass $m$ at height $h$ is $E_p=mgh$; take $g=9{,}81\,\text{m}\!\cdot\!\text{s}^{-2}$.
\end{indication}

\begin{solution}
\textbf{1.} Heating the ice to its melting temperature requires
\[
Q_1=mc_{\text{ice}}\Delta T
=1{,}0\times2{,}1\times10^3\times100
=2{,}1\times10^5\,\text{J}.
\]

\textbf{2.} During melting, the temperature remains at $0\,{}^\circ\text{C}$:
\[
Q_2=mL_f=1{,}0\times3{,}34\times10^5
=3{,}34\times10^5\,\text{J}.
\]

\textbf{3.} To heat the liquid water from $0$ to $100\,{}^\circ\text{C}$,
\[
Q_3=mc_{\text{water}}\Delta T
=1{,}0\times4{,}18\times10^3\times100
=4{,}18\times10^5\,\text{J}.
\]

\textbf{4.} Complete vaporization requires a further
\[
Q_4=mL_v=1{,}0\times2{,}26\times10^6
=2{,}26\times10^6\,\text{J}.
\]
This stage alone therefore requires a few megajoules.

\textbf{5.} Vaporization is by far the most energy-intensive stage:
\[
Q_4=2{,}26\times10^6\,\text{J}
>Q_3=4{,}18\times10^5\,\text{J}
>Q_2=3{,}34\times10^5\,\text{J}
>Q_1=2{,}1\times10^5\,\text{J}.
\]
In particular, vaporization requires about $(2{,}26\times10^6)/(4{,}18\times10^5)\approx5{,}4$ times as much energy as heating the liquid water from $0$ to $100\,{}^\circ\text{C}$.

For a mass $m$ falling through $h=10\,\text{m}$, $E_p=mgh$. Setting $mgh=Q_3$ gives
\[
m=\frac{Q_3}{gh}
=\frac{4{,}18\times10^5}{9{,}81\times10}
\approx4{,}26\times10^3\,\text{kg}.
\]
A mass of about $4{,}3$ tonnes would therefore have to fall ten metres to release as much energy as is needed to heat one kilogram of water from $0$ to $100\,{}^\circ\text{C}$! This is enormous and explains why water heating accounts for a significant share of domestic energy use. Water also has a very large specific heat capacity, especially compared with common metals: it is about ten times that of copper, for example (see the following exercises).

\textbf{6.} Condensation is the reverse of vaporization: the vapour releases the latent heat
\[
Q_{\text{rel}}=m_vL_v.
\]
For $m_v=1{,}0\times10^{-2}\,\text{kg}$,
\[
Q_{\text{rel}}=1{,}0\times10^{-2}\times2{,}45\times10^6
=2{,}45\times10^4\,\text{J}.
\]
Thus, even a small mass of condensing vapour can release a non-negligible amount of heat, received by the windscreen and its surroundings.
\end{solution}
\end{exo}

\begin{exo}[Evapotranspiration from a large tree: a natural air conditioner?]{odg-evapotranspiration-arbre}
\seoready{true}
\keywords{order of magnitude, evapotranspiration, tree, latent heat, cooling power, air conditioner, COP}

On a hot, dry day, a large tree with an adequate water supply may evapotranspire about $500\,\text{L}=5{,}0\times10^{-1}\,\text{m}^3$ of water. Since transpiration occurs mainly in daylight, assume that it is spread over twelve hours. The tree canopy covers a ground area $A_{\text{tree}}=50\,\text{m}^2$. The density of water and its latent heat of vaporization at atmospheric pressure are
\[
\rho_{\text{water}}=1{,}0\times10^3\,\text{kg}\!\cdot\!\text{m}^{-3},
\qquad L_v=2{,}45\times10^6\,\text{J}\!\cdot\!\text{kg}^{-1}.
\]

\begin{enumerate}
  \item Calculate the mass of water evapotranspired during the day and the energy required to vaporize it.
  \item Explain why this produces a cooling effect.
  \item Calculate the tree's average cooling power per square metre during the twelve hours of active transpiration.
  \item For comparison, consider an air conditioner supplied with electrical power $P_{\mathrm{elec}}=2{,}0\times10^3\,\text{W}$, with cooling coefficient of performance $\mathrm{COP}=3{,}0$, cooling a room of area $A_{\mathrm{room}}=20\,\text{m}^2$. By definition, $\mathrm{COP}=P_{\mathrm{cool}}/P_{\mathrm{elec}}$. Calculate its cooling power and cooling power per square metre, then compare them with the tree's values.
  \item Why can we not conclude that the tree and the air conditioner provide exactly the same perceived cooling, even if their powers per unit area have the same order of magnitude?
  \item Could many houseplants cool a home? (General knowledge: the answer depends on biological processes.)
\end{enumerate}

\begin{indication}
Use $m=\rho V$, $Q_{\mathrm{evap}}=mL_v$, and $P=Q/\tau$. For the air conditioner, $P_{\mathrm{cool}}=\mathrm{COP}\,P_{\mathrm{elec}}$.
\end{indication}

\begin{solution}
\textbf{1.} The corresponding mass is
\[
m=\rho_{\text{water}}V
=1{,}0\times10^3\times5{,}0\times10^{-1}
=5{,}0\times10^2\,\text{kg}.
\]
The vaporization energy is
\[
Q_{\mathrm{evap}}=mL_v
=5{,}0\times10^2\times2{,}45\times10^6
=1{,}225\times10^9\,\text{J},
\]
whose order of magnitude is one gigajoule.

\textbf{2.} Vaporization is endothermic. It draws this energy from the leaves, soil, and surrounding air, thereby lowering the temperature of the leaves and their immediate environment.

\textbf{3.} Twelve hours correspond to $\tau=12\times3600=4{,}32\times10^4\,\text{s}$. Hence
\[
P_{\text{tree}}=\frac{Q_{\mathrm{evap}}}{\tau}
=\frac{1{,}225\times10^9}{4{,}32\times10^4}
\approx2{,}84\times10^4\,\text{W},
\]
and
\[
\frac{P_{\text{tree}}}{A_{\text{tree}}}
=\frac{2{,}84\times10^4}{50}
\approx5{,}7\times10^2\,\text{W}\!\cdot\!\text{m}^{-2}.
\]

\textbf{4.} The useful cooling power is
\[
P_{\text{cool}}=\mathrm{COP}\,P_{\mathrm{elec}}
=3{,}0\times2{,}0\times10^3
=6{,}0\times10^3\,\text{W}.
\]
Per unit floor area,
\[
\frac{P_{\text{cool}}}{A_{\text{room}}}
=\frac{6{,}0\times10^3}{20}
=3{,}0\times10^2\,\text{W}\!\cdot\!\text{m}^{-2}.
\]
In this model, the tree's evapotranspiration power per unit area is almost twice that of the air conditioner.

\textbf{5.} This comparison concerns energy fluxes only. An air conditioner cools a controlled enclosed volume, whereas the vapour from a tree disperses into the atmosphere and wind brings in warm air. Perceived cooling depends on ventilation, humidity, ambient temperature, and distance from the tree. Transpiration also varies with sunlight, wind, air humidity, species, and available water, and can decrease sharply during drought. The tree additionally provides shade, which directly reduces the solar radiation received by the ground and by people. The calculation therefore compares orders of magnitude, not equal thermal comfort.

\textbf{6.} In practice, houseplants provide very little cooling. In most plants, light encourages stomata to open, allowing water vapour to escape. Indoor light is generally weak, so stomata open less and transpiration decreases. In a poorly ventilated room, rising humidity also progressively slows evaporation.

Many plants may therefore cause slight local evaporative cooling, but cannot replace an air conditioner. Removing heat durably would require moist air to be vented outdoors. Intense artificial lighting might stimulate transpiration, but its electrical energy would itself end up almost entirely as heat in the room. CAM plants adapted to arid environments, including cacti, are an exception: their stomata open mainly at night.
\end{solution}
\end{exo}

\begin{exo}[Mixing two masses of water]{calorimetrie-melange-eau}
\seoready{true}
\keywords{calorimetry, mixing, heat capacity, specific heat capacity, Joseph Black, equilibrium temperature}

A mass $m_1=0{,}200\,\text{kg}$ of water at $T_1=80\,^\circ\text{C}$ is poured into a vessel already containing $m_2=0{,}300\,\text{kg}$ of water at $T_2=15\,^\circ\text{C}$. The vessel is an ideal calorimeter: heat loss to the surroundings and the heat capacity of the vessel itself are neglected. Recall that $Q=mc\Delta T$, where $c=4{,}18\times10^3\,\text{J}\!\cdot\!\text{kg}^{-1}\!\cdot\!\text{K}^{-1}$ is the specific heat capacity of liquid water.

\begin{enumerate}
  \item Justify that the heat released by the hot water is entirely received by the cold water. Write the conservation equation involving $m_1$, $m_2$, $c$, $T_1$, $T_2$, and the equilibrium temperature $T_{eq}$.
  \item Deduce an expression for $T_{eq}$ and calculate its numerical value.
  \item Calculate the heat $Q$ released by the hot water during mixing.
  \item Can mixtures of the same substance be used to determine $c$? Explain.
\end{enumerate}

\begin{indication}
Because the calorimeter is isolated and rigid, the total internal energy $U_1+U_2$ is conserved: the heat released by one mass is received by the other, with the opposite sign.
\end{indication}

\begin{solution}
\textbf{1.} Conservation of the total internal energy of the system gives
\[
m_1c(T_{eq}-T_1)+m_2c(T_{eq}-T_2)=0.
\]

\textbf{2.} The common coefficient $c$ cancels, giving
\[
T_{eq}=\frac{m_1T_1+m_2T_2}{m_1+m_2}
=\frac{0{,}200\times80+0{,}300\times15}{0{,}200+0{,}300}
=41\,^\circ\text{C}.
\]

\textbf{3.} The heat released by the hot water is
\[
Q=m_1c(T_1-T_{eq})
=0{,}200\times4{,}18\times10^3\times(80-41)
\approx3{,}26\times10^4\,\text{J}.
\]
The cold water receives the same quantity:
\[
m_2c(T_{eq}-T_2)
=0{,}300\times4{,}18\times10^3\times26
\approx3{,}26\times10^4\,\text{J}.
\]

\textbf{4.} No. For two masses of the same substance, $c$ multiplies both terms and cancels. The measured equilibrium temperature is therefore independent of $c$ and is simply the mass-weighted mean of the initial temperatures. Measuring a specific heat capacity by the method of mixtures requires another substance whose heat capacity is known, as in the next exercise.
\end{solution}
\end{exo}

\begin{exo}[Determining the specific heat capacity of a metal by the method of mixtures]{calorimetrie-methode-melanges}
\seoready{true}
\keywords{calorimetry, method of mixtures, specific heat capacity, calorimeter, water equivalent}

As Joseph Black did in the eighteenth century, we wish to measure the specific heat capacity $c_m$ of an unknown metal by the \emph{method of mixtures}. A sample of mass $m=0{,}150\,\text{kg}$ is heated to $T_m=95\,^\circ\text{C}$ and quickly immersed in a calorimeter containing $m_e=0{,}250\,\text{kg}$ of water, with $c_e=4{,}18\times10^3\,\text{J}\!\cdot\!\text{kg}^{-1}\!\cdot\!\text{K}^{-1}$, initially at $T_e=18\,^\circ\text{C}$. The measured equilibrium temperature is $T_{eq}=22\,^\circ\text{C}$. Initially, neglect the heat capacity of the calorimeter vessel, stirrer, and thermometer.

\begin{enumerate}
  \item Write the energy-conservation equation for the isolated system consisting of the metal and water, with $c_m$ as the only unknown.
  \item Deduce an expression for $c_m$ and calculate it. Which common metal does it approximately match? Reference values, in $\text{J}\!\cdot\!\text{kg}^{-1}\!\cdot\!\text{K}^{-1}$, are: aluminium $9{,}0\times10^2$, iron $4{,}5\times10^2$, copper $3{,}9\times10^2$, and lead $1{,}3\times10^2$.
  \item The calorimeter vessel actually absorbs some heat. Model this by a \emph{water equivalent} $\mu=1{,}5\times10^{-2}\,\text{kg}$, as though the calorimeter were an additional mass $\mu$ of water. Recalculate $c_m$ and comment on the direction of the correction.
  \item Why must the sample be immersed quickly and the calorimeter be well insulated from the surrounding air?
\end{enumerate}

\begin{indication}
The metal releases $Q_m=mc_m(T_m-T_{eq})$, received by the water and, in Question 3, by the calorimeter: $Q_m=(m_e+\mu)c_e(T_{eq}-T_e)$.
\end{indication}

\begin{solution}
\textbf{1.} Energy conservation gives
\[
mc_m(T_m-T_{eq})=m_ec_e(T_{eq}-T_e).
\]

\textbf{2.} Thus
\[
c_m=\frac{m_ec_e(T_{eq}-T_e)}{m(T_m-T_{eq})}
=\frac{0{,}250\times4{,}18\times10^3\times(22-18)}{0{,}150\times(95-22)}
\approx3{,}82\times10^2\,\text{J}\!\cdot\!\text{kg}^{-1}\!\cdot\!\text{K}^{-1}.
\]
This is very close to the value for copper.

\textbf{3.} Including the calorimeter's water equivalent,
\[
c_m=\frac{(m_e+\mu)c_e(T_{eq}-T_e)}{m(T_m-T_{eq})}
=\frac{(0{,}250+0{,}015)\times4{,}18\times10^3\times4}{0{,}150\times73}
\approx4{,}05\times10^2\,\text{J}\!\cdot\!\text{kg}^{-1}\!\cdot\!\text{K}^{-1}.
\]
The correction increases $c_m$: neglecting the vessel underestimates the heat actually released by the metal and therefore underestimates its specific heat capacity.

\textbf{4.} The method assumes an isolated system throughout the measurement. Rapid immersion limits the time during which the hot sample can exchange heat with the surrounding air before reaching the water, while good insulation limits losses to the surroundings as equilibrium is established. Any unaccounted heat transfer would distort the balance in Question 1 and hence the inferred value of $c_m$. This is precisely the kind of experimental care already taken by Black, and later by Lavoisier and Laplace with their ice calorimeter.
\end{solution}
\end{exo}

\begin{exo}[Food calories and metabolic power]{calorimetrie-calories-alimentaires}
\seoready{true}
\keywords{calorie, kilocalorie, food Calorie, mechanical equivalent, metabolic power, units}

A nutrition label indicates that an adult consumes about $2000\,\text{kcal}$ per day on average. The kilocalorie is a non-SI unit still used in nutrition; its SI conversion is $1\,\text{kcal}=4{,}18\times10^3\,\text{J}$.

\begin{enumerate}
  \item Convert this daily intake into joules.
  \item Deduce the adult's average metabolic power in watts, assuming the energy is consumed uniformly over 24 hours.
  \item An energy bar is labelled ``$210\,\text{kcal}$''. What mass of water, initially at $20\,^\circ\text{C}$, could be heated to boiling point, $100\,^\circ\text{C}$, with the same energy if all of it were converted into heat? Take $c_{\text{water}}=4{,}18\times10^3\,\text{J}\!\cdot\!\text{kg}^{-1}\!\cdot\!\text{K}^{-1}$.
  \item In what forms is this energy actually used by the human body? Give a qualitative answer.
\end{enumerate}

\begin{indication}
For Question 2, use $\mathcal P=E/\Delta t$. For Question 3, first convert the labelled energy into joules, then solve $Q=mc\Delta T$ for $m$.
\end{indication}

\begin{solution}
\textbf{1.} In SI units,
\[
E=2000\times4{,}18\times10^3
=8{,}36\times10^6\,\text{J}.
\]

\textbf{2.} The average power is
\[
\mathcal P=\frac{8{,}36\times10^6}{8{,}64\times10^4}\approx97\,\text{W}.
\]
Over a whole day, a person therefore uses about as much average power as an incandescent light bulb. This order of magnitude is often surprising and illustrates the value of converting food calories into familiar physical units.

\textbf{3.} The energy of the bar is
\[
Q=210\times4{,}18\times10^3\approx8{,}78\times10^5\,\text{J}.
\]
With $\Delta T=80\,\text{K}$,
\[
m=\frac{Q}{c\Delta T}
=\frac{8{,}78\times10^5}{4{,}18\times10^3\times80}
\approx2{,}63\,\text{kg}.
\]
Thus, to order of magnitude, one energy bar contains enough energy to bring about $2{,}5\,\text{kg}$ of tap water to boiling point.

\textbf{4.} The body does not consume energy in the sense of making it disappear; it transforms chemical energy stored in nutrients. Some supports muscular work, organ function, cellular electrical gradients, molecular transport, and chemical synthesis. Some may be stored chemically as glycogen or fat. Ultimately, most of the energy used is dissipated as heat, helping maintain body temperature before being transferred to the environment. A small fraction also leaves the body as chemical energy in waste.
\end{solution}
\end{exo}
